% 900 characters, including spaces, required
% - Highlights are a short collection of bullet point statements (3-5) that concisely convey to the reader the recent advances in the area, including emerging concepts and/or distinctions, that constitute a main motivation for the discussion developed in the article. 
% - As Highlights focus on recent developments, conclusions and future directions should instead be discussed in the Concluding Remarks section and/or the Outstanding Questions box.
% - The Highlights are not called out in the text.
% - The Highlights do not count towards the total number of allowed display elements in the manuscript (limit of 6).
% - When submitting the manuscript files via Editorial Manager, please upload Highlights as a separate word file using the designated heading.
 
\begin{framed}
  \section*{Highlights}
  \begin{itemize}
  \item Deep networks and techniques from deep reinforcement learning
    have greatly widened the scope of computational simulations of language
    emergence in communities of interactive agents.
  \item Thanks to these modern tools, language emergence can now be
    studied among agents that receive realistic perceptual input, must
    solve complex tasks cooperatively or competitively, and can engage
    in flexible multi-turn verbal and non-verbal interactions.
  \item With great simulation power comes great need for new analysis
    methods: a budding area of research focuses on understanding
    the general characteristics of the deep agents' emergent language.
  \item Another line of research wants to deliver on the promise of
    interactive AI, exploring the functional role of emergent
    language in improving machine-machine and human-machine communication.
  \end{itemize}
\end{framed}
